% PSMFlows pipeline, explained simply.
% Companion to main.tex (the formal write-up): this document trades rigor for clarity.
% Build: latexmk -pdf pipeline.tex
\documentclass[11pt]{article}
\usepackage[margin=1.1in]{geometry}
\usepackage{amsmath,amssymb}
\usepackage{algorithm}
\usepackage{algpseudocode}
\usepackage{booktabs}
\usepackage[colorlinks=true,linkcolor=blue!60!black]{hyperref}

\newcommand{\D}{\mathcal{D}}
\newcommand{\E}{\mathbb{E}}
\newcommand{\R}{\mathbb{R}}

\title{The PSMFlows Pipeline, in Plain RL Terms}
\author{}
\date{Companion to \texttt{main.tex} --- August 2026}

\begin{document}
\maketitle

\section{The problem in one paragraph}

We are given a fixed dataset $\D=\{(s_i,a_i,s'_i)\}$ of transitions collected by some
unknown behavior policy, \emph{with no rewards}. Later, someone hands us a reward
function, and we must produce a good policy \emph{without any further environment
interaction} (zero-shot offline RL). Methods like FB and PSM solve this by learning
successor measures --- ``if I act like policy $\pi$ from state $s$, which states will I
visit?'' --- for a whole family of policies at once. Their weakness: during training they
evaluate candidate \emph{actions} that may never appear in the data. The value estimates
for such out-of-distribution (OOD) actions are unconstrained, the argmax hunts for
exactly those over-estimated actions, and the resulting policy is garbage precisely where
the data is thin. PSMFlows removes this failure mode by construction: \textbf{every
action the method ever considers is produced by a generative model of the data itself.}

\section{The key idea: index policies by flow noise}

Train a conditional flow $G_\theta(s,u)$ that maps a noise vector
$u\sim\mathcal{N}(0,I_{d_a})$ to an action that looks like the dataset's actions at state
$s$ (behavior cloning with a flow, exactly the BC part of FQL). Because the noise
dimension equals the action dimension, $G_\theta(s,\cdot)$ is an invertible map: every
noise vector decodes to a data-like action, and every action has a unique noise vector
that produces it.

Now the trick: \emph{hold the noise fixed}. The deterministic policy
$\pi_u(s)=G_\theta(s,u)$ takes, at every state, ``the kind of action the data takes at
$s$, in direction $u$.'' Different $u$'s give different, reproducible behaviors (our D2
diagnostic measures exactly this), so the noise vector $u$ becomes a \emph{policy index}.
Optimizing over $u\in\R^{d_a}$ replaces optimizing over raw actions --- and any $u$ of
typical norm decodes to an action the data supports. The hard constraint ``stay on the
data manifold'' becomes the easy constraint ``keep $\|u\|$ typical.''

\section{The three training stages}

\subsection{Stage A --- learn the behavior flow (reward-free BC)}

Standard conditional flow matching: sample a dataset pair $(s,a)$, a noise $u_0$, a time
$t$, and regress the velocity network $v_\theta(s,\,(1-t)u_0+ta,\,t)$ onto $(a-u_0)$.
Integrating $v_\theta$ from $u$ over $t\in[0,1]$ gives the decoder $G_\theta(s,u)$; a
one-step student network is distilled from the integrated flow for fast acting. In the
code this is \texttt{agent=fql agent.bc\_only=true} (the critic is never trained; rewards
are never read). \textbf{The flow is then frozen forever.} Freezing matters: the policy
family that the next stages build on is stationary, unlike FB's actor which moves during
training.

\subsection{Stage B --- invert the flow over the dataset (preimages)}

Stage C will need to know, for every dataset transition, \emph{which latent the behavior
``used''}: the $u$ with $G_\theta(s,u)\approx a$. Since $G_\theta(s,\cdot)$ is
invertible, run its ODE backwards from $a$ to get the point preimage
$u^\star=E_\theta(s,a)$. Because a whole basin of latents decodes to nearly the same
action, we also fit a distribution over that basin, the relaxed posterior
\[
  q_\alpha(u\mid s,a)\;\propto\;
  \mathcal{N}(u;0,I)\,\exp\!\bigl(-\alpha\,\|G_\theta(s,u)-a\|\bigr),
\]
by importance-sampled EM started from a Laplace (inverse-Jacobian) guess. The
temperature $\alpha$ sets how close a decoded action must be to $a$ to count; the
effective sample size (ESS) of the EM tells us whether the fit is trustworthy (D3 gates
on it). Everything is precomputed once (\texttt{tools/precompute\_preimages.py}) and
stored, one posterior per transition, next to a sidecar recording exactly which flow
checkpoint produced it.

Sanity property worth remembering: if the flow perfectly modeled the behavior
distribution, inverted latents would be exactly $\mathcal{N}(0,I)$ distributed, so
$\|u^\star\|^2$ must look $\chi^2_{d_a}$. That is the D3 typicality test --- a direct
check that train-time latents and test-time latents come from the same distribution.

\subsection{Stage C --- successor measures over the latent-indexed family}

A successor measure $M^\pi(s,a,X)$ is the discounted time a policy spends in the state
set $X$ after taking $a$ at $s$. PSMFlows learns them for the whole family
$\{\pi_{u'}\}$ at once, in factored form
\[
  M^{u\to u'}(s,X)\;\approx\;\psi(s,u',u)^\top\varphi(X),
\]
read as: ``take the action indexed by $u$ now, then follow policy $\pi_{u'}$ forever.''
$\varphi$ is a state basis (kept near-orthonormal by a regularizer), $\psi$ carries all
policy dependence. The training signal is the standard measure-matching TD loss (as
PSM/FB): for a batch transition $(s,a,s')$,
\begin{itemize}
  \item the \emph{action slot} $u$ is a sample from the stored preimage posterior
        $q_\alpha(\cdot\mid s,a)$ --- the dataset action, expressed as a latent;
  \item the \emph{policy index} $u'$ is a fresh draw (half $\mathcal{N}(0,I)$, half
        permuted behavior preimages, clipped to the typical set);
  \item the bootstrap target is $\psi(s',u',u')$: at the next state, the continuation
        policy \emph{is} the indexed policy, the same $(Q,V)$ relationship as ordinary
        TD.
\end{itemize}
Note what is absent: no actor network, no argmax over actions, no decoded action anywhere
in training. The flow only decodes at \emph{acting} time. This is the sense in which
every action the backup implicitly evaluates is data-supported.

\section{Inference: reward in, policy out}

Given rewards $r$ on (a relabeled sample of) the dataset, two cheap steps:
\begin{enumerate}
  \item \textbf{Task inference.} Because $\varphi$ is near-orthonormal,
        $w=\E[\,r(s')\,\varphi(s')\,]$ is the least-squares projection of the reward onto
        the basis, and $\psi(s,u',u)^\top w \approx Q$-value of ``do $u$, then follow
        $\pi_{u'}$'' for this reward. Closed form, no training.
  \item \textbf{Flow-GPI.} At each state draw $K$ candidate action-latents and $M$ policy
        indices, score every pair by $\psi(s,u',u)^\top w$ (minus a small ensemble-
        disagreement penalty), take the best candidate $u$, and decode it through the
        frozen flow: $a=G_\theta(s,u)$. This is generalized policy improvement over the
        latent family: at least as good as the best indexed policy in the sampled set,
        while every candidate action stays a flow decode.
\end{enumerate}

\section{The full algorithm}

\begin{algorithm}[h]
\caption{PSMFlows (as implemented on \texttt{feat/psm-integration})}
\begin{algorithmic}[1]
\State \textbf{Input:} reward-free dataset $\D$; later, a reward signal $r$
\Statex \emph{--- Stage A: behavior flow (frozen afterwards) ---}
\For{each BC step}
  \State sample $(s,a)\sim\D$, $u_0\sim\mathcal{N}(0,I)$, $t\sim U[0,1]$
  \State regress $v_\theta(s,(1-t)u_0+ta,\,t)\to(a-u_0)$;\quad distill one-step decoder
\EndFor
\Statex \emph{--- Stage B: preimages (precomputed once) ---}
\For{each $(s_i,a_i)\in\D$}
  \State $u_i^\star \gets E_\theta(s_i,a_i)$ by backward ODE; Jacobian $J_i$ by autodiff
  \State fit $q_\alpha(u\mid s_i,a_i)$ by IS-EM from the Laplace guess; record ESS
\EndFor
\State validate: round-trip error, $\chi^2$ typicality of $\{u_i^\star\}$, mean ESS (D3 gate)
\Statex \emph{--- Stage C: successor-measure training ---}
\For{each TD step, batch $\{(s_i,a_i,s'_i)\}$}
  \State $u_i\sim q_\alpha(\cdot\mid s_i,a_i)$ \Comment{dataset action as a latent}
  \State $u'_i\gets$ mix of $\mathcal{N}(0,I)$ and permuted $\{u_j\}$, clipped to $\pm u_{\text{clip}}$
  \State $M_{ij}\gets\psi(s_i,u'_i,u_i)^\top\varphi(s'_j)$;\quad
         $\bar M_{ij}\gets\bar\psi(s'_i,u'_i,u'_i)^\top\bar\varphi(s'_j)$ \Comment{targets}
  \State descend the measure TD loss $+\;\lambda_\perp\|\tfrac1B\sum_j\varphi\varphi^\top-I\|_F^2$;
         Polyak-update targets
\EndFor
\Statex \emph{--- Inference (zero-shot, per reward) ---}
\State $w\gets\frac1N\sum_i r(s'_i)\,\varphi(s'_i)$, normalized
\For{each encountered state $s$}
  \State draw $u_{1:K},\,u'_{1:M}\sim\mathcal{N}(0,I)$ clipped;\quad
         $\hat u\gets\arg\max_k \max_m \psi(s,u'_m,u_k)^\top w$
  \State act $a\gets G_\theta(s,\hat u)$ \Comment{one-step decode; always data-like}
\EndFor
\end{algorithmic}
\end{algorithm}

\section{The diagnostics, and what each one certifies}

\begin{center}
\small
\begin{tabular}{@{}llp{5.9cm}l@{}}
\toprule
 & Tool & Question, in plain terms & Gate \\
\midrule
D1 & \texttt{validate\_flow\_fidelity} & does the flow reproduce the data's actions? & beat random control \\
D2 & \texttt{eval\_fixed\_u\_rollouts} & does a fixed $u$ mean a distinct, repeatable policy? & ratio $<1$ \\
D3 & \texttt{validate\_flow\_inversion} & can we trust the preimages? & ESS, typicality, roundtrip \\
D4 & \texttt{latent\_q\_sanity} & with oracle rewards, does GPI actually control? & $\ge50\%$ of FQL \\
\bottomrule
\end{tabular}
\end{center}

\noindent
Each stage only starts once the previous stage's gate passes, so a failure is localized:
a bad D1 means the flow, a bad D3 means the inversion, a bad D4 with good D1--D3 means
the representation or GPI --- never a mystery spread across the whole pipeline.

\end{document}
